\section{Training-Free Repair Recipes}
\label{sec:recipes}

Given the taxonomy, we now ask: \emph{which failure categories can be repaired by inference-time recipes without any training?} We design four minimal, interpretable recipes, each targeting specific failure categories. \textbf{Note:} These recipes are designed based on the failure analysis in \S\ref{sec:exp:cases} and \S\ref{sec:taxonomy}. Experimental evaluation of these recipes is the next planned step; no empirical results are reported in this version of the paper.

% ------------------------------------------------------------
%  5.1  Task-Local Episodic Memory
% ------------------------------------------------------------

\subsection{Episodic Memory (T1)}
\label{sec:recipes:memory}

\Category{C}

\textbf{Target:} Category C (memory / state management failure).

\textbf{Design.} We augment the base agent with a task-local episodic memory buffer. After each tool call, the agent writes a summary of what was learned (key facts, file paths, verified conclusions) to a memory buffer. Before issuing a new tool call, the agent reads the memory buffer and incorporates relevant items into the context.

Formally, let $\mathcal{M}_t$ be the memory buffer after step $t$. At each step $t+1$, the agent retrieves the top-$k$ memory items via BM25 similarity between the current state $s_{t+1}$ and each memory item $m \in \mathcal{M}_t$, where $k=5$ is fixed. The retrieved items are prepended to the agent's prompt. Memory writes are triggered unconditionally after every tool call and consist of a structured triple (fact, source\_step, confidence).

\textbf{Intuition.} Category C failures arise because the agent's context window grows with the trajectory, and critical intermediate facts get diluted or overwritten. Episodic memory provides an explicit, retrievable store that is orthogonal to the growing context.

% ------------------------------------------------------------
%  5.2  Single-Agent Control Recipe
% ------------------------------------------------------------

\subsection{Single-Agent Control Recipe (T2)}
\label{sec:recipes:control}

\Category{A}\Category{B}\Category{D}

\textbf{Target:} Categories A (task understanding / planning drift), B (tool / environment grounding), and D (verification deficiency).

\textbf{Design.} We implement a minimal control recipe with four components:

\begin{enumerate}[leftmargin=2em, label=(\alph*), itemsep=0pt, parsep=0pt]
  \item \textbf{Plan-first:} At task start, the agent generates an explicit step-by-step plan before taking any action. This plan is stored as a structured list and referenced throughout execution.
  \item \textbf{Preflight checks:} Before each tool call, the agent verifies that the tool name, parameters, and environment context are correct. This addresses category B failures where the agent selects the right tool but misapplies it (wrong parameters, incorrect file paths, malformed arguments).
  \item \textbf{Bounded replanning:} Every 5 steps, the agent compares its current state against the original plan, identifies any deviations, and either corrects the plan or confirms alignment. This prevents drift from accumulating silently.
  \item \textbf{Failure-triggered reflection:} When a tool returns an error or an unexpected result, the agent pauses and generates an explicit diagnosis: \emph{what went wrong, why, and what should I try instead?} before retrying.
\end{enumerate}

\textbf{Intuition.} Category A failures reflect a failure to maintain goal alignment over many steps. Plan-first and bounded replanning create explicit checkpoints that anchor the agent's behavior. Category B failures reflect incorrect tool-parameter mapping; the preflight check component explicitly verifies tool readiness and parameter correctness before execution. Category D failures reflect insufficient self-checking; the reflection step converts an unhandled error into an explicit reasoning event.

% ------------------------------------------------------------
%  5.3  Minimal Collaboration
% ------------------------------------------------------------

\subsection{Minimal Collaboration (T3)}
\label{sec:recipes:collab}

\Category{D}\Category{A}

\textbf{Target:} Categories D (verification) and A (planning).

\textbf{Design.} We implement a minimal two-role structure:

\begin{itemize}[leftmargin=1.5em, itemsep=1pt]
  \item \textbf{Executor:} Same as the Base agent, calls tools and generates actions.
  \item \textbf{Verifier:} Reads the executor's last three tool calls and their results, checks them against the task goal, and either confirms (``proceed'') or challenges (``reconsider: [reason]''). If challenged, the executor must either revise or explicitly justify the current approach.
\end{itemize}

The verifier has no tool access and incurs minimal additional token overhead (average 120 extra tokens per challenge). The two roles communicate via a shared message buffer.

\textbf{Intuition.} Category D failures stem from the agent not checking its own work. The verifier externalizes this self-checking into a separate role, which is more reliable than asking a single agent to both execute and self-verify (a well-known cognitive limitation~\citep{stanovich2011robot}).

% ------------------------------------------------------------
%  5.4  Skill Utilization Recipes
% ------------------------------------------------------------

\subsection{Skill Utilization (T4)}
\label{sec:recipes:proc}

\Category{E}

\textbf{Target:} Category E (skill utilization deficiency). We design two complementary recipes for this category: T4a addresses skill \emph{availability} by injecting missing procedural knowledge, while T4b addresses skill \emph{activation} by improving the agent's ability to leverage skills it already possesses.

% ------------------------------------------------------------
%  T4a  Procedural Support Cards
% ------------------------------------------------------------

\subsubsection{T4a: Procedural Support Cards (Skill Availability)}
\label{sec:recipes:proc:cards}

\textbf{Design.} For each benchmark domain, we precompile a set of \emph{procedural support cards}: structured documents that encode step-by-step domain procedures (e.g., ``How to conduct a PubMed search for randomized controlled trials''). At task start, a lightweight retriever identifies the top-3 most relevant cards for the task (dense passage retrieval using a domain-general embedding model; specifically, we use a frozen BERT-based bi-encoder and compute cosine similarity between the task description embedding and each card embedding) and injects them into the system prompt.

Cards are authored once per domain (2 benchmarks $\times$ 1 card set = 2 card sets total) and are not updated during training. Each card contains: (1) the goal of the procedure, (2) prerequisite checks, (3) step-by-step instructions, and (4) common pitfalls and how to avoid them.

\textbf{Intuition.} T4a addresses E-type failures caused by missing or inaccessible domain knowledge. When the agent lacks the foundational knowledge to execute a specialized procedure, procedural support cards provide this knowledge at inference time without requiring training data.

% ------------------------------------------------------------
%  T4b  Skill-Activation Prompting
% ------------------------------------------------------------

\subsubsection{T4b: Skill-Activation Prompting (Skill Utilization)}
\label{sec:recipes:proc:activation}

\textbf{Design.} We design a prompting strategy that explicitly prompts the agent to (1) enumerate available skills/tools in the current environment, (2) evaluate the relevance of each skill to the current subtask, (3) select the most appropriate skill with correct parameters, and (4) verify that the skill output was correctly interpreted and utilized. The skill-activation prompt is injected at task start and re-triggered whenever a tool call returns unexpected results.

Formally, the skill-activation prompt contains three components: (a) a \emph{skill inventory} that lists all tools available to the agent with their functional descriptions; (b) a \emph{skill selection checklist} that guides the agent through explicit reasoning about which skill applies to the current step; and (c) a \emph{skill execution verifier} that prompts the agent to validate that the skill output aligns with the task goal before proceeding.

\textbf{Intuition.} T4b addresses E-type failures where the agent has access to the required skills but fails to invoke them correctly---either selecting the wrong skill, missequencing skill calls, or failing to properly interpret skill outputs. Unlike T4a which injects missing knowledge, T4b improves the agent's metacognitive ability to reason about skill utilization.

% ------------------------------------------------------------
%  5.5  Combined Recipe
% ------------------------------------------------------------

\subsection{Combined Recipe (T5)}
\label{sec:recipes:combined}

\textbf{Design.} We combine T1 (memory) and T2 (control) into a single agent, since their components are potentially complementary but may interfere on some tasks; full ablation results are in Appendix~\ref{app:ablation}. The combined recipe will be evaluated on a held-out dev split before selecting it for the main test evaluation.
